% ============================================================
\section{Software Ecosystem and Research Infrastructure}
\label{sec:software}
% ============================================================

The scientific value of any measurement system depends not only on hardware precision
but on the transparency of the software stack and the researcher's unrestricted
access to raw signal data.
Proprietary black-box systems that withhold signal-level access and expose only
opaque aggregated metrics do not qualify for rigorous academic study.
The FlexTail ecosystem was designed from the outset for complex longitudinal
data analysis.

\subsection{Configurable acquisition rates}

To accommodate experimental designs with very different resolution requirements, the
system supports a continuously adjustable sampling rate.

For epidemiological long-term recordings over multiple days --- characterising daily
sitting time, postural load distribution in a nursing shift, or nightly sleep
patterns --- a rate of 1--25\,Hz provides the necessary temporal resolution while
conserving battery life.

For high-dynamic biomechanical investigations --- gait analysis, sit-to-stand
transitions, resistance training kinematics, or rehabilitation assessments requiring
sub-degree temporal resolution of rapid movements --- the rate can be raised to
100\,Hz \cite{walkling2025}.

\subsection{The \emph{fiffi} middleware layer}

Bluetooth Low Energy (\acs{BLE}) connectivity dropouts represent a persistent source of
data loss in wearable field studies.
The \emph{fiffi} middleware, written in C and running directly on the sensor
microcontroller, eliminates this failure mode through autonomous buffering.

When the \acs{BLE} link to the gateway device is interrupted --- because the participant
leaves their smartphone at a desk, moves through a shielded area, or enters
electromagnetic interference --- \emph{fiffi} immediately transitions the sensor to
standalone data-logger mode.
All angular and \acs{IMU} data are written to internal flash memory at the configured
sampling rate.
On reconnection, \emph{fiffi} initiates a prioritised, lossless flash-synchronisation
transfer, restoring the complete continuous time series without gaps \cite{walkling2025}.

This architecture ensures that multi-day field recordings are not compromised by
the connectivity conditions of uncontrolled environments such as hospital wards,
construction sites, or participants' homes.

\subsection{Raw data access and export formats}

At each time step, the system delivers an unfiltered data packet containing:

\begin{enumerate}
	\item Sacral Euler angles (pitch, roll) from the base \acs{IMU} (\si{\degree}).
	\item 18 local sagittal flexion angles from the strain gauges (\si{\degree}).
	\item 18 local torsion angles around the longitudinal axis (\si{\degree}).
	\item Global 3D Cartesian coordinates for all 18 segments (mm), computed via the
	      \ac{DH} chain.
	\item Pre-computed aggregate metrics: lumbar partial angle, thoracic partial angle,
	      total trunk torsion.
\end{enumerate}

Researchers visualise, annotate, and export these datasets through the
\textbf{FlexLab Studio} desktop application.
FlexLab Studio provides a 3D real-time rendering of the reconstructed spinal shape
alongside time-series plots, event annotation tools, and batch export functionality.

For maximum read/write performance with large longitudinal datasets, the software
provides the proprietary \acf{RSF} (a compressed binary format optimised for the
multi-channel time series structure) and the columnar Apache Parquet format.
For compatibility with established statistical environments (SPSS, R, Python/Pandas,
MATLAB), standard CSV export is also available.

\subsection{The \emph{flexlib} Python library}

Post-hoc analysis of large longitudinal datasets --- which may contain hundreds of
millions of individual time points from multi-week recordings --- is facilitated by
\emph{flexlib}, an open-architecture Python library.
\emph{Flexlib} loads proprietary \acs{RSF} data directly into Python and exposes
validated algorithms for bending detection, ergonomic scoring, and postural
state classification.

The computationally intensive \ac{DH} transformation pipeline is implemented as a
Rust back-end compiled into the Python package via Maturin, providing near-native
execution speed for bulk processing without requiring participants to perform
computations on the sensor hardware.

\subsection{Hysteresis-based postural state detection}

High-frequency sensor time series from the human body contain inherent noise from
minor muscle fasciculations, respiratory mechanics, and micro-vibrations.
A naive threshold algorithm --- ``lumbar angle above $X$\si{\degree} defines the
state \emph{flexed}'' --- produces rapid false toggling (flickering) around the
boundary whenever the signal oscillates near the threshold due to this noise.

\emph{Flexlib} addresses this with Schmitt-trigger-style hysteresis filters: each
postural state is associated with a separate entry threshold and a lower exit
threshold.
The system transitions into the state only when the signal exceeds the entry
threshold, and leaves it only when the signal drops below the (lower) exit
threshold.
The hysteresis band effectively suppresses noise-driven toggling and produces
stable, physiologically interpretable state sequences \cite{walkling2025}.

\begin{figure}[ht]
	\centering
	\begin{tikzpicture}
		\begin{axis}[
				name=signal,
				width=0.82\textwidth,
				height=4.5cm,
				xmin=0, xmax=10,
				ymin=10, ymax=45,
				xlabel={},
				ylabel={Lumbar angle (°)},
				ylabel style={font=\small},
				yticklabel style={font=\small},
				xticklabels={},
				tick style={draw=gray!50},
				axis line style={draw=gray!60},
				grid=major,
				grid style={dashed, gray!25},
				clip=false,
			]
			% Simulated noisy lumbar angle signal
			\addplot[minktecBlue, thick, smooth] coordinates {
					(0,20)(0.5,22)(1,28)(1.5,31)(2,33)(2.5,34)(3,34.5)(3.5,33)(4,30)(4.5,27)
					(5,25)(5.5,24)(6,23)(6.5,22)(7,20)(7.5,19)(8,18)(8.5,20)(9,23)(9.5,24)(10,25)
				};
			% Entry threshold (upper)
			\draw[red!70!black, dashed, thick] (axis cs:0,32) -- (axis cs:10,32)
			node[right, font=\footnotesize, red!70!black] {entry ($\theta_\text{on}$)};
			% Exit threshold (lower)
			\draw[orange!80!black, dashed, thick] (axis cs:0,26) -- (axis cs:10,26)
			node[right, font=\footnotesize, orange!80!black] {exit ($\theta_\text{off}$)};
			% Highlight active zone
			\fill[minktecBlue!15, opacity=0.5] (axis cs:1.8,26) rectangle (axis cs:4.3,32);
		\end{axis}

		% State output panel below
		\begin{axis}[
				name=state,
				at={(signal.below south west)}, anchor=above north west,
				yshift=-4mm,
				width=0.82\textwidth,
				height=2.0cm,
				xmin=0, xmax=10,
				ymin=-0.1, ymax=1.3,
				xlabel={Time (arbitrary units)},
				xlabel style={font=\small},
				ylabel={State},
				ylabel style={font=\small},
				yticklabel style={font=\small},
				xticklabel style={font=\small},
				ytick={0,1},
				yticklabels={off, on},
				tick style={draw=gray!50},
				axis line style={draw=gray!60},
				grid=major,
				grid style={dashed, gray!25},
			]
			% Step function: off until signal exceeds entry, on until below exit
			\addplot[minktecNavy, very thick, const plot] coordinates {
					(0,0)(1.8,0)(1.8,1)(4.3,1)(4.3,0)(10,0)
				};
		\end{axis}
	\end{tikzpicture}
	\caption{Schmitt-trigger hysteresis in the postural state detector. The sensor signal (top) crosses the entry threshold~$\theta_\text{on}$ at $t \approx 1.8$, triggering the \emph{flexed} state (bottom). The state persists until the signal drops below the lower exit threshold~$\theta_\text{off}$ at $t \approx 4.3$, preventing false toggling from noise oscillations near a single boundary.}
	\label{fig:hysteresis}
\end{figure}

\subsection{Real-time interaction via \emph{fiffi\_unleashed}}

For experimental designs that require acute behavioural modification, the
\emph{fiffi\_unleashed} API provides bidirectional real-time access to the sensor.

On the input side, the complete 3D data stream is available at the configured
sampling rate and can be fed directly into external simulation environments such as
Unity or OpenSim for real-time skeletal animation or biomechanical modelling.

On the output side, the API exposes the onboard vibration motor.
Researchers can trigger specific vibration patterns programmatically --- for example,
activating a brief pulse whenever the cumulative time in lumbar hyperkyphosis exceeds
a configurable threshold --- to provide unobtrusive tactile nudges that prompt
postural correction without interrupting the participant's task \cite{walkling2025}.
